\chapter{结果讨论示例章}
\label{chap:discussion}

本章展示结果讨论部分的常见格式，包括结果对比、机理分析、工程应用建议等。

\section{结果对比分析}
\subsection{不同方法对比}
本研究结果与已有研究成果的对比见表\ref{tab:result_comparison}。

\begin{table}[htbp]
    \caption{本研究结果与已有成果对比}
    \label{tab:result_comparison}
    \centering
    \begin{tabular}{lcccc}
        \toprule
        研究来源 & 方法类型 & 主要参数 & 结果数值 & 相对误差 \\
        \midrule
        本研究 & 数值模拟 & $E=210$ GPa & 385.6 MPa & --- \\
        文献\cite{smith2022} & 试验研究 & $E=205$ GPa & 380.2 MPa & 1.4\% \\
        文献\cite{zhang2020} & 理论分析 & $E=210$ GPa & 392.5 MPa & 1.8\% \\
        \bottomrule
    \end{tabular}
\end{table}

从对比结果可以看出，本研究结果与已有研究成果吻合良好，验证了研究方法的可靠性。

\subsection{误差分析}
误差来源主要包括以下几个方面：
\begin{enumerate}
    \item 数值离散化误差：有限元网格尺寸引入的离散化误差；
    \item 材料参数误差：材料本构参数与实际值存在偏差；
    \item 边界条件误差：简化边界条件与实际工况的差异。
\end{enumerate}

\section{机理分析}
\subsection{力学机理}
通过数值分析和试验验证，揭示了以下力学机理：

（1）荷载传递机理：荷载通过材料内部传递，应力分布呈现明显的不均匀性。

（2）变形机理：材料变形包括弹性变形和塑性变形两个阶段，塑性变形集中于高应力区域。

（3）破坏机理：材料破坏起源于应力集中区域，裂纹扩展导致最终破坏。

\subsection{影响因素分析}
影响研究结果的关键因素包括：
\begin{itemize}
    \item 材料因素：材料的力学性能直接影响研究对象的响应特性。
    \item 几何因素：结构几何形状决定了应力分布和变形模式。
    \item 边界因素：边界条件对数值分析结果有显著影响。
\end{itemize}

\section{工程应用建议}
\subsection{设计建议}
基于研究结果，提出以下设计建议：
\begin{enumerate}
    \item 参数选择建议：建议设计参数控制在合理范围内。
    \item 尺寸优化建议：建议结构尺寸比例控制在特定范围内以提高效率。
    \item 材料选择建议：建议选用性能稳定的材料。
\end{enumerate}

\subsection{施工建议}
施工过程中应注意以下事项：
\begin{itemize}
    \item 控制施工误差，确保几何尺寸满足设计要求；
    \item 严格控制材料质量，确保材料性能达标；
    \item 施工顺序应按照设计方案执行。
\end{itemize}

\section{研究局限性}
本研究存在以下局限性：

（1）研究范围局限：本研究仅针对特定工况进行研究，其他工况需进一步验证。

（2）方法局限：数值分析方法存在一定的简化假设，可能影响结果的准确性。

（3）数据局限：试验数据有限，可能影响结论的普适性。